\chapter{Planteamiento del Problema}

Los herbarios como cualquier banco de datos debe contar con herramientas confiables que faciliten la gestión de información al igual que la centralización, auditoria y mantenimiento de las mismas. Esta tarea puede representar un desafió enorme dada la cantidad de muestras y de información presente en ellas. La única aproximación viable a un problema así es la construcción de un gestor de información digital que permita automatizar las labores de centralización y curaduria de un cuerpo de datos tan extenso como el que puede tener una institución como esta. el diseño he implementación va a requerir la aplicación de diferentes tecnologías y frameworks orientados al aprendizaje automático y sistemas distribuidos , en concreto la principal situación problema a la que queremos dar solución es :

Es posible generar un modelo de aprendizaje profundo y visión por computador el cual sea capaz de centralizar y gestionar las diferentes fuentes de información presentes como son muestras de ejemplares y de esta manera tener una herramienta versátil para la curaduria de estas instituciones 

\clearpage
%\newpage
%\chapter{Hipótesis}